% !TEX program = xelatex
\documentclass[12pt,a4paper]{article}

% ══════════════════════════════════════════════════════════════════════
% PACKAGES
% ══════════════════════════════════════════════════════════════════════
\usepackage{fontspec}
\usepackage{xeCJK}
\setCJKmainfont{SimSun}[BoldFont=SimHei,ItalicFont=KaiTi]
\usepackage[english]{babel}
\usepackage{amsmath,amssymb,amsthm}
\usepackage{mathtools}
\usepackage{bm}
\usepackage{graphicx}
\usepackage{hyperref}
\usepackage{geometry}
\usepackage{enumitem}
\usepackage{booktabs}
\usepackage{tikz-cd}
\usepackage{xcolor}
\usepackage{cleveref}
\usepackage{bbm}
\usepackage{array}
\usepackage{multirow}
\usepackage{makecell}
\usepackage{longtable}
\usepackage{float}

\geometry{margin=2.5cm}

\hypersetup{
    colorlinks=true,
    citecolor=blue,
    urlcolor=blue,
    pdftitle={Topological Signatures of Adoption Phase Transitions in Multi-Expert Systems},
    pdfauthor={SCX},
    pdfsubject={SCX Topological Adoption},
}

% ══════════════════════════════════════════════════════════════════════
% THEOREM ENVIRONMENTS
% ══════════════════════════════════════════════════════════════════════
\theoremstyle{definition}
\newtheorem{definition}{Definition}[section]
\newtheorem{theorem}{Theorem}[section]
\newtheorem{proposition}{Proposition}[section]
\newtheorem{corollary}{Corollary}[section]
\newtheorem{lemma}{Lemma}[section]
\newtheorem{conjecture}{Conjecture}[section]
\newtheorem{problem}{Open Problem}[section]

\theoremstyle{remark}
\newtheorem{remark}{Remark}[section]
\newtheorem{example}{Example}[section]

\theoremstyle{definition}
\newtheorem*{definition*}{Definition}
\newtheorem*{theorem*}{Theorem}
\newtheorem*{proposition*}{Proposition}
\newtheorem*{corollary*}{Corollary}
\newtheorem*{lemma*}{Lemma}
\newtheorem*{remark*}{Remark}
\newtheorem*{problem*}{Open Problem}

% ══════════════════════════════════════════════════════════════════════
% CUSTOM COMMANDS
% ══════════════════════════════════════════════════════════════════════
\newcommand{\R}{\mathbb{R}}
\newcommand{\C}{\mathbb{C}}
\newcommand{\Z}{\mathbb{Z}}
\newcommand{\N}{\mathbb{N}}
\newcommand{\Q}{\mathbb{Q}}
\newcommand{\F}{\mathbb{F}}
\newcommand{\E}{\mathbb{E}}
\newcommand{\I}{\mathbb{I}}

\newcommand{\Od}{O(d)}
\newcommand{\SOd}{SO(d)}

% SCX-specific
\newcommand{\CEC}{\mathcal{E}}
\newcommand{\SCX}{\textsc{SCX}}
\newcommand{\Yajie}{\textsc{Yajie}}
\newcommand{\Spring}{\textsc{Spring}}
\newcommand{\Cercis}{\textsc{Cercis}}
\newcommand{\CECstar}{|\CEC|^*}
\newcommand{\betaone}{\beta_1}
\newcommand{\Deltabeta}{\Delta\beta_1}

% ── Operators ──
\DeclareMathOperator{\Hom}{Hom}
\DeclareMathOperator{\im}{im}
\DeclareMathOperator{\kerOp}{ker}
\DeclareMathOperator{\tr}{tr}
\DeclareMathOperator{\vol}{vol}
\DeclareMathOperator{\spec}{Spec}

\begin{document}

% ══════════════════════════════════════════════════════════════════════
% TITLE PAGE
% ══════════════════════════════════════════════════════════════════════
\title{\textbf{Topological Signatures of Adoption Phase Transitions}\\[6pt]
       \large in Multi-Expert Audit Systems}

\author{SCX Research Group\\[4pt]
        \texttt{scx@nousresearch.com}}

\date{July 2026\\[4pt]
      \small\textit{Prototype --- v0.1}}

\maketitle

\begin{abstract}
We establish a topological classification of adoption phase transitions in multi-expert audit systems by linking the gauge-theoretic first Betti number $\beta_1 = (M-1)(M-2)/2$ of the \SCX\ expert graph to the six-stage lock-in trajectory $\{S_1,\ldots,S_6\}$ of protocol adoption dynamics. The invariant $\beta_1$, which counts independent harmonic 1-forms---equivalently, incommensurable flat gauge connections---exhibits discontinuous jumps at each increment of the expert count $M$: $0 \to 1 \to 3 \to 6 \to 10 \to \cdots$ We propose that these topological phase transitions induce corresponding adoption phase transitions, where each jump in $\beta_1$ reduces the threshold Cumulative Evidentiary Corpus size $\CECstar$ required for lock-in and enables the system to traverse successive stages $S_1 \to S_2 \to \cdots \to S_6$. We derive a scaling relationship between $\beta_1$ and the critical CEC value, compute the topological fingerprint of each adoption stage, and identify $M=4$ ($\beta_1=3$) as the minimal expert count for robust adoption dynamics. A conjecture is formulated positing that $\CECstar \propto \beta_1^{-1}$ in the large-$M$ limit, implying that topological complexity accelerates adoption convergence.
\end{abstract}

\tableofcontents
\newpage

% ══════════════════════════════════════════════════════════════════════
% 1. INTRODUCTION
% ══════════════════════════════════════════════════════════════════════
\section{Introduction}

\subsection{Motivation}

The \SCX\ framework for multi-expert data quality auditing rests on two mathematically rigorous pillars. The first is gauge-theoretic: the \SCX\ expert graph, formalized as a 2-dimensional cell complex $K_{\SCX}$, possesses a discrete Hodge decomposition whose harmonic subspace encodes the fundamentally irreducible inconsistencies among expert assessments~\cite{scx_gauge_formalized,scx_fiber_bundle}. The dimension of this harmonic subspace is the first Betti number
\begin{equation}\label{eq:beta_def}
  \beta_1(M) = \frac{(M-1)(M-2)}{2},
\end{equation}
which depends solely on the number of experts $M$ and is independent of the number of configurations $N$.

The second pillar is game-theoretic: the \Yajie\ Protocol's adoption dynamics are modeled as a lock-in Markov process over six stages $\mathcal{S} = \{S_1,\ldots,S_6\}$, from emergence to universal adoption (``天下大同'')~\cite{yajie_protocol}. The transition from $S_2$ (Critical Mass) to $S_3$ (Moat Formation/NPE) is governed by the cumulative evidentiary corpus $\CEC$ reaching a critical threshold $\CECstar$, at which the non-proliferation equilibrium becomes the unique Nash equilibrium of the adoption game.

These two pillars---topological and game-theoretic---have been developed independently. Yet they share a common parameter: the expert count $M$. This paper proposes that the independent development is a historical artifact, not a structural necessity. We demonstrate that the topological invariant $\beta_1(M)$ provides a \emph{signature} that predicts the qualitative character of adoption phase transitions, and that the discontinuous jumps in $\beta_1$ as $M$ increases map systematically onto the lock-in stages $S_1 \to S_6$.

\subsection{Core Insight}

The core insight is simple but consequential. The first Betti number $\beta_1$ counts the number of \emph{independent harmonic modes} in the \SCX\ expert graph---equivalently, the number of incommensurable flat gauge connections, or the dimension of the moduli space of irreducible inconsistencies. When $\beta_1 = 0$ ($M=2$), there are \emph{no} harmonic modes: any inconsistency between two experts is, by the Hodge decomposition, necessarily a gradient (exact form) and can be resolved by a global gauge transformation. Detection of systematic inconsistencies reduces to trivial pairwise comparison.

When $\beta_1 = 1$ ($M=3$), the \emph{first} harmonic mode appears. For the first time, the system admits an inconsistency that is \emph{not} reducible to a gradient---a genuinely topological obstruction to global gauge alignment. This is the mathematical origin of the ``irreducible inconsistency'' that the \SCX\ framework's multi-expert consensus mechanism is designed to detect.

As $M$ grows, $\beta_1$ jumps discontinuously: $M=2 \to \beta_1=0$, $M=3 \to \beta_1=1$, $M=4 \to \beta_1=3$, $M=5 \to \beta_1=6$, $M=10 \to \beta_1=36$, and so on. Each jump represents a qualitative change in the topological complexity of the inconsistency space. We propose that these jumps are the topological correlates of the adoption phase transitions identified in the game-theoretic analysis.

\subsection{Structure of the Paper}

Section~2 reviews the gauge-theoretic foundation: the \SCX\ 2-complex, the discrete Hodge decomposition, and the derivation of $\beta_1(M)$. Section~3 reviews the six-stage adoption lock-in model and the critical CEC threshold $\CECstar$. Section~4 establishes the central relationship between $\beta_1$ and $\CECstar$, deriving a scaling law and identifying $M=4$ as a topological threshold. Section~5 maps the $\beta_1$ jumps onto the adoption stages $S_1$--$S_6$ and proposes a topological phase diagram. Section~6 discusses implications for protocol design: minimum expert count, optimal auditor pool size, and the role of topological redundancy. Section~7 states our main conjecture, and Section~8 concludes with open problems.

% ══════════════════════════════════════════════════════════════════════
% 2. GAUGE-THEORETIC FOUNDATION
% ══════════════════════════════════════════════════════════════════════
\section{Gauge-Theoretic Foundation: The \texorpdfstring{$\beta_1$}{β₁} Invariant}

\subsection{The SCX Expert Graph as a 2-Complex}

We begin by recalling the construction of the \SCX\ expert graph as a 2-dimensional cell complex $K_{\SCX}$, following~\cite{scx_gauge_formalized}.

Let there be $M$ experts and $N$ configurations (parameter settings). The vertices are the pairs $(k,m)$ with $k \in \{1,\ldots,N\}$ and $m \in \{1,\ldots,M\}$, giving $|V| = NM$. The edges are of two types:
\begin{itemize}[nosep]
  \item \textbf{Parameter edges}: $(k,m) \to (k+1,m)$, connecting the same expert across adjacent configurations. There are $M(N-1)$ such edges.
  \item \textbf{Expert edges}: $(k,i) \to (k,j)$, connecting different experts at the same configuration. There are $NM(M-1)$ such edges.
\end{itemize}
The total edge count is $|E| = NM^2 - M$.

The faces are elementary quadrilaterals $(k,i) \to (k,j) \to (k+1,j) \to (k+1,i)$, of which there are $|F| = (N-1)M(M-1)/2$. The resulting 2-complex carries a discrete Hodge theory with boundary operators $d_0: \R^{|V|} \to \R^{|E|}$ (incidence matrix) and $d_1: \R^{|E|} \to \R^{|F|}$ (curl operator).

\begin{theorem}[Homotopy Type of $K_{\SCX}$, Theorem 3.1 of~\cite{scx_gauge_formalized}]\label{thm:homotopy}
  The \SCX\ 2-complex is homotopy equivalent to the complete graph on $M$ vertices:
  \begin{equation}
    K_{\SCX} \simeq K_M.
  \end{equation}
  Consequently, its first homology group is free of rank
  \begin{equation}\label{eq:beta_explicit}
    \beta_1(K_{\SCX}) = \beta_1(K_M) = \binom{M-1}{2} = \frac{(M-1)(M-2)}{2}.
  \end{equation}
\end{theorem}

\begin{corollary}[Betti Number Independence]\label{cor:betti_indep}
  $\beta_1$ depends only on the expert count $M$ and is independent of the configuration count $N$.
\end{corollary}

The independence of $N$ is a deep structural fact: the topological complexity of the inconsistency space is determined entirely by the number of experts, not by the volume of data they process. This is the first hint that $M$ is the controlling parameter for qualitative adoption dynamics.

\subsection{Harmonic Modes and Irreducible Inconsistency}

By the discrete Hodge decomposition~\cite{scx_fiber_bundle},
\begin{equation}\label{eq:hodge}
  \Omega^1(K_{\SCX}) = \im(d_0) \oplus \kerOp(\Delta_1) \oplus \im(d_1^\dagger),
\end{equation}
where $\Omega^1$ is the space of edge assignments (the gauge potential $A$), $\im(d_0)$ is the space of pure-gauge (exact) configurations, $\kerOp(\Delta_1)$ is the space of harmonic 1-forms, and $\im(d_1^\dagger)$ is the space of co-exact 1-forms.

The harmonic subspace $\kerOp(\Delta_1)$ has dimension $\beta_1$. Each basis harmonic 1-form $h_i$ ($i = 1,\ldots,\beta_1$) satisfies
\begin{equation}
  d_0^\dagger h_i = 0 \quad\text{(divergence-free)},\qquad
  d_1 h_i = 0 \quad\text{(curl-free)},
\end{equation}
meaning it is simultaneously impossible to express as a gradient and impossible to detect by measuring circulation around elementary faces.

This is the mathematical formalization of \emph{irreducible inconsistency}: a harmonic mode represents a pattern of expert disagreement that cannot be eliminated by any global gauge transformation (re-calibration) and cannot be localized to any specific face (pairwise comparison). It is a genuinely topological obstruction to consensus.

\begin{theorem}[Harmonic Moduli Space, Theorem 3.4 of~\cite{scx_gauge_formalized}]\label{thm:moduli}
  The moduli space of flat gauge connections on $K_{\SCX}$ has dimension
  \begin{equation}
    \dim \mathcal{M}_{\text{flat}} = \beta_1 \cdot \dim(\mathfrak{g}),
  \end{equation}
  where $\mathfrak{g}$ is the Lie algebra of the gauge group. For $O(d)$ gauge,
  $\dim(\mathcal{M}_{\text{flat}}) = \beta_1 \cdot d(d-1)/2$;
  for Abelian gauge ($\R^d$), $\dim(\mathcal{M}_{\text{flat}}) = \beta_1 \cdot d$.
\end{theorem}

\subsection{The $\beta_1$ Sequence}

We compute $\beta_1$ for representative values of $M$:

\begin{table}[H]
  \centering
  \caption{First Betti number $\beta_1(M)$ and $\dim\mathcal{M}_{\text{flat}}$ for $d=3$}
  \label{tab:beta_sequence}
  \begin{tabular}{@{}c c c c c@{}}
    \toprule
    $M$ & $\beta_1$ & $\Delta\beta_1$ (jump) & $\dim\mathcal{M}_{\text{flat}}$ (Abelian) & $\dim\mathcal{M}_{\text{flat}}$ ($O(3)$) \\
    \midrule
    2 & 0  & --- & 0  & 0 \\
    3 & 1  & $+1$ & 3  & 3 \\
    4 & 3  & $+2$ & 9  & 9 \\
    5 & 6  & $+3$ & 18 & 18 \\
    6 & 10 & $+4$ & 30 & 30 \\
    7 & 15 & $+5$ & 45 & 45 \\
    8 & 21 & $+6$ & 63 & 63 \\
    9 & 28 & $+7$ & 84 & 84 \\
    10& 36 & $+8$ & 108& 108 \\
    \bottomrule
  \end{tabular}
\end{table}

Three observations are immediate:
\begin{enumerate}[nosep]
  \item $\beta_1$ grows quadratically: $\beta_1 \sim M^2/2$ as $M \to \infty$.
  \item The jumps $\Delta\beta_1(M) = \beta_1(M+1) - \beta_1(M) = M-1$ grow linearly, meaning each additional expert adds \emph{more} independent harmonic modes than the previous one.
  \item The first non-zero $\beta_1$ occurs at $M=3$ ($\beta_1=1$), and the jump from $M=3$ to $M=4$ adds $\Delta\beta_1 = 2$, tripling the total ($\beta_1: 1 \to 3$).
\end{enumerate}

The topological interpretation is that $M=2$ represents a ``tree-level'' system with no cycles, $M=3$ introduces exactly one independent cycle (a triangle), and $M=4$ introduces three independent cycles (a tetrahedron). Each cycle is a potential locus of irreducible inconsistency.

\subsection{Detection Capacity via $\beta_1$}

The $\beta_1$ harmonic modes provide \emph{independent detection channels} for systemic inconsistency. By the Hoeffding concentration argument that underpins the \SCX\ noise detection guarantee (Theorem~1 of~\cite{scx_theory}), if the system has $\beta_1$ independent harmonic modes, each providing an independent signal of global inconsistency, then the effective number of independent ``auditors'' is not $M$ but
\begin{equation}\label{eq:M_eff}
  M_{\text{eff}} = M \cdot (1 + c \cdot \beta_1),
\end{equation}
where $c > 0$ is a coupling constant representing the per-mode detection efficiency. The intuition is that each harmonic mode provides an additional dimension along which expert inconsistency can be detected, amplifying the effective statistical power of the multi-expert ensemble.

\emph{Justification.} Each harmonic mode $h_i \in \kerOp(\Delta_1)$ defines a linear functional $\langle A, h_i\rangle$ on the gauge potential. If $A$ contains a non-zero harmonic component, at least one such functional is non-zero, and the Hoeffding bound on detection probability acquires an additional factor from the independent test along this harmonic direction. Since the $\beta_1$ modes are mutually orthogonal, the tests are independent, and the effective sample size increases multiplicatively.

% ══════════════════════════════════════════════════════════════════════
% 3. ADOPTION PHASE TRANSITIONS
% ══════════════════════════════════════════════════════════════════════
\section{Adoption Phase Transitions: The \texorpdfstring{$S_1$--$S_6$}{S1--S6} Model}

\subsection{The Lock-In Markov Process}

The adoption dynamics of the \Yajie\ Protocol are modeled as a Markov process over six stages~\cite{yajie_protocol}:
\begin{equation}
  \mathcal{S} = \{S_1, S_2, S_3, S_4, S_5, S_6\},
\end{equation}
with an upper-triangular transition matrix (irreversible lock-in):
\begin{equation}
  p_{ij} = 0 \quad \text{for all } j < i.
\end{equation}

The six stages are:

\begin{description}[nosep,leftmargin=2em]
  \item[$S_1$ --- Emergence (零散萌芽).] Protocol launched, $\CEC = \varnothing$. Accuracy is low. Adoption driven by early adopters. No lock-in.
  \item[$S_2$ --- Critical Mass (临界质量).] Cumulative audit volume reaches $\CECstar$. Accuracy advantage becomes statistically significant. Lock-in begins to emerge but remains contestable.
  \item[$S_3$ --- Moat Formation / NPE (护城河形成).] Non-Proliferation Equilibrium becomes the unique Nash equilibrium. Protocol becomes the default standard. Ecosystem externalities accelerate.
  \item[$S_4$ --- Institutionalization (制度化).] Protocol embedded in legal and institutional fabric. Switching costs become prohibitive---not because of technical superiority, but because the institutional environment has been built around the protocol.
  \item[$S_5$ --- Epistemic Closure (认知闭合).] CEC becomes so comprehensive that no meaningful challenge to its authority is possible. The definition of data quality \emph{is} what the protocol measures. Lock-in is complete.
  \item[$S_6$ --- Universal Harmony (天下大同).] CEC internationalized. Protocol governed by multilateral framework. Adoption universal and irreversible. True absorbing state.
\end{description}

\subsection{The Critical CEC Threshold $\CECstar$}

The transition from $S_2$ to $S_3$ is governed by the \textbf{CEC Critical Value Theorem} (Theorem~2 of~\cite{yajie_protocol}):

\begin{theorem}[CEC Critical Value, adapted from~\cite{yajie_protocol}]\label{thm:cec_critical}
  There exists a unique critical CEC size $\CECstar$ such that, for all $|\CEC| > \CECstar$, the all-adopt strategy profile $(A,\ldots,A)$ is the unique Nash equilibrium of the non-proliferation game. The critical value satisfies
  \begin{equation}\label{eq:cec_star}
    \Delta(\CECstar) = -\lambda,
  \end{equation}
  where $\Delta(|\CEC|) = \theta(|\CEC|) - \theta_0$ is the accuracy advantage of the incumbent over a zero-CEC entrant, and $\lambda > 0$ is the fragmentation cost parameter.
\end{theorem}

The accuracy advantage $\Delta$ is a monotonically increasing function of $|\CEC|$, typically modeled as
\begin{equation}\label{eq:theta_cec}
  \theta(|\CEC|) = \theta_\infty - (\theta_\infty - \theta_0) e^{-\gamma |\CEC|},
\end{equation}
where $\theta_\infty$ is the asymptotic accuracy ceiling, $\theta_0$ is initial accuracy, and $\gamma$ is the learning rate.

\subsection{The Role of Expert Count $M$}

The accuracy function $\theta(|\CEC|)$ implicitly depends on $M$ through the multi-expert detection mechanism. As established in \SCX\ Theorem~1~\cite{scx_theory}, the false-positive rate of multi-expert consensus detection decays as $\exp(-2M\Delta^2)$, and the effective expert count $M_{\text{eff}} = M/(1+(M-1)\bar{\rho})$ accounts for expert correlation $\bar{\rho}$.

In the topological framework developed here, we propose that the effective detection rate is enhanced by the harmonic mode count $\beta_1(M)$, so that
\begin{equation}\label{eq:theta_enhanced}
  \theta(|\CEC|; M) = \theta_\infty - (\theta_\infty - \theta_0) e^{-\gamma \cdot |\CEC| \cdot (1 + \alpha \cdot \beta_1(M))},
\end{equation}
where $\alpha \in (0,1)$ is a topological coupling constant representing the per-mode detection efficiency gain.

The consequence is that larger $M$ (and hence larger $\beta_1$) accelerates the approach to the accuracy ceiling $\theta_\infty$, reducing the CEC size required to reach the critical threshold $\CECstar$.

% ══════════════════════════════════════════════════════════════════════
% 4. THE $\beta_1$--$\CECstar$ RELATIONSHIP
% ══════════════════════════════════════════════════════════════════════
\section{The \texorpdfstring{$\beta_1$--$\CECstar$}{β₁--|ℰ|*} Relationship}

\subsection{Derivation of the Critical Threshold}

From equation~\eqref{eq:theta_enhanced} and the critical condition $\Delta(\CECstar) = -\lambda$ (Theorem~\ref{thm:cec_critical}), we solve for $\CECstar$:

\begin{align}
  \theta(\CECstar; M) - \theta_0 &= -\lambda \notag\\
  \theta_\infty - (\theta_\infty - \theta_0) e^{-\gamma \CECstar (1 + \alpha\beta_1)} - \theta_0 &= -\lambda \notag\\
  (\theta_\infty - \theta_0)(1 - e^{-\gamma \CECstar (1 + \alpha\beta_1)}) &= -\lambda \notag\\
  e^{-\gamma \CECstar (1 + \alpha\beta_1)} &= 1 + \frac{\lambda}{\theta_\infty - \theta_0}.
\end{align}

Let $\kappa = \lambda/(\theta_\infty - \theta_0) > 0$ be the normalized fragmentation cost. Then
\begin{align}
  -\gamma \CECstar (1 + \alpha\beta_1) &= \ln(1 + \kappa) \notag\\
  \CECstar(M) &= -\frac{\ln(1+\kappa)}{\gamma} \cdot \frac{1}{1 + \alpha\beta_1(M)}.
\end{align}

\begin{proposition}[Topological Reduction of Critical CEC]\label{prop:cec_reduction}
  The critical CEC size as a function of the expert count is
  \begin{equation}\label{eq:cec_star_formula}
    \CECstar(M) = \frac{\CECstar_0}{1 + \alpha \cdot \beta_1(M)},
  \end{equation}
  where $\CECstar_0 = -\ln(1+\kappa)/\gamma$ is the critical CEC size in the absence of topological enhancement ($\beta_1 = 0$, i.e., $M=2$).
\end{proposition}

This is the central result of the paper. It states that the topological complexity $\beta_1(M)$ \emph{directly reduces} the CEC accumulation required to achieve adoption lock-in. The mechanism is the multiplicative enhancement of detection efficiency through independent harmonic modes.

\begin{corollary}[Topological Acceleration Factor]\label{cor:acceleration}
  The relative reduction in required CEC is
  \begin{equation}
    \frac{\CECstar_0 - \CECstar(M)}{\CECstar_0} = \frac{\alpha\beta_1(M)}{1 + \alpha\beta_1(M)}.
  \end{equation}
  For large $M$, $\beta_1 \sim M^2/2$, and the reduction approaches unity: $\CECstar(M) \to 0$ as $M \to \infty$.
\end{corollary}

\subsection{Numerical Evaluation}

We evaluate $\CECstar(M)$ for representative values of $M$, normalized against the baseline $\CECstar_0$ ($M=2$, $\beta_1=0$):

\begin{table}[H]
  \centering
  \caption{Normalized critical CEC $\CECstar(M)/\CECstar_0$ for $\alpha = 0.1$ and $\alpha = 0.3$}
  \label{tab:cec_reduction}
  \begin{tabular}{@{}c c c c c@{}}
    \toprule
    $M$ & $\beta_1$ & $1+\alpha\beta_1$ ($\alpha=0.1$) & $\CECstar/\CECstar_0$ ($\alpha=0.1$) & $\CECstar/\CECstar_0$ ($\alpha=0.3$) \\
    \midrule
    2 & 0  & 1.0 & 1.000 & 1.000 \\
    3 & 1  & 1.1 & 0.909 & 0.769 \\
    4 & 3  & 1.3 & 0.769 & 0.526 \\
    5 & 6  & 1.6 & 0.625 & 0.357 \\
    6 & 10 & 2.0 & 0.500 & 0.250 \\
    7 & 15 & 2.5 & 0.400 & 0.182 \\
    8 & 21 & 3.1 & 0.323 & 0.137 \\
    9 & 28 & 3.8 & 0.263 & 0.106 \\
    10& 36 & 4.6 & 0.217 & 0.085 \\
    \bottomrule
  \end{tabular}
\end{table}

The qualitative pattern is clear: $\CECstar$ drops sharply from $M=2$ to $M=3$ (the first harmonic mode appears), again from $M=3$ to $M=4$ (the jump $\beta_1: 1 \to 3$), and continues to decay as $M$ grows. Each unit increase in $M$ after $M=4$ yields diminishing relative returns, consistent with the $\sim 1/M^2$ asymptotic scaling of $1/(1+\alpha\beta_1)$.

\subsection{The $M=4$ Threshold}

A striking feature of the data is the magnitude of the jump from $M=3$ to $M=4$:
\begin{itemize}[nosep]
  \item $M=3$: $\beta_1 = 1$, $\CECstar/\CECstar_0 = 0.909$ (only 9\% reduction).
  \item $M=4$: $\beta_1 = 3$, $\CECstar/\CECstar_0 = 0.769$ (23\% reduction).
\end{itemize}
The $\beta_1$ jump of $+2$ (tripling from 1 to 3) produces a disproportionate effect because it is the first transition in which the number of harmonic modes exceeds one: at $M=3$, there is a single inconsistency channel; at $M=4$, there are three, enabling triangulation of inconsistencies. We therefore identify

\begin{proposition}[Topological Robustness Threshold]\label{prop:M4_threshold}
  $M=4$ ($\beta_1=3$) is the minimal expert count for \emph{topologically robust} adoption dynamics, defined as the point where $\CECstar$ has been reduced by at least 20\% from baseline and where the system possesses at least two independent inconsistency detection channels beyond the first ($\beta_1 \geq 3$).
\end{proposition}

This is reminiscent of the $K \geq 3$ minimum in the protocol governance paper~\cite{protocol_governance}, where three maintainers are the minimum for consensus deviation detection. The topological analysis provides independent geometric grounds for a similar bound, shifted by one: $M=4$ rather than $M=3$, because the topological requirement is not merely majority voting but \emph{triangulation of inconsistency} in a multidimensional harmonic space.

% ══════════════════════════════════════════════════════════════════════
% 5. TOPOLOGICAL PHASE DIAGRAM
% ══════════════════════════════════════════════════════════════════════
\section{Topological Phase Diagram of Adoption}

\subsection{Mapping $\beta_1$ Jumps to Adoption Stages}

We propose the following correspondence between topological phase transitions ($\beta_1$ jumps) and adoption phase transitions ($S_i \to S_j$):

\begin{table}[H]
  \centering
  \caption{Topological signatures of adoption phase transitions}
  \label{tab:phase_mapping}
  \begin{tabular}{@{}c c c c p{4.5cm}@{}}
    \toprule
    $\beta_1$ Jump & $M$ range & Adoption Transition & Topological Mechanism & Physical Interpretation \\
    \midrule
    $0 \to 1$ & $2 \to 3$ & $S_1 \to S_2$ & First harmonic mode appears & System acquires capacity to detect irreducible inconsistency\\
    $1 \to 3$ & $3 \to 4$ & $S_2 \to S_3$ & Triangulation in harmonic space & NPE becomes unique equilibrium; $|\CEC| > \CECstar$ achievable\\
    $3 \to 6$ & $4 \to 5$ & $S_3 \to S_4$ & Rich harmonic topology & Institutional lock-in; harmonic modes span all inconsistency types\\
    $6 \to 10$ & $5 \to 6$ & $S_4 \to S_5$ & Dense harmonic coverage & Epistemic closure; no inconsistency escapes detection\\
    $10 \to 36$ & $6 \to 10$ & $S_5 \to S_6$ & Asymptotic saturation & Universal harmony; $\CECstar \to 0$, adoption universal\\
    \bottomrule
  \end{tabular}
\end{table}

The mapping is not one-to-one in a strict mathematical sense---adoption transitions depend on additional economic and institutional parameters beyond the topological. But the mapping is \emph{structurally predictive}: each $\beta_1$ jump corresponds to a \emph{necessary condition} for the adoption transition to become possible. Without sufficient topological complexity (harmonic modes), the detection capacity is insufficient for the CEC to reach $\CECstar$ within any reasonable time frame.

\subsection{Phase Diagram}

Figure~\ref{fig:phase_diagram} (conceptual) illustrates the topological phase diagram in the $(M, \CECstar)$ plane.

\begin{table}[H]
  \centering
  \caption{Phase regions in $(M, \beta_1)$ space}
  \label{tab:phase_regions}
  \begin{tabular}{@{}c c c c c@{}}
    \toprule
    Region & $M$ & $\beta_1$ & Adoption Stage & Detection Character \\
    \midrule
    \textbf{Trivial} & $2$ & $0$ & Pre-$S_1$ & No harmonic channels; pairwise comparison only \\
    \textbf{Emergent} & $3$ & $1$ & $S_1$--$S_2$ & One inconsistency mode; fragile detection\\
    \textbf{Robust} & $4$--$5$ & $3$--$6$ & $S_2$--$S_3$ & Triangulation possible; strong detection\\
    \textbf{Institutional} & $6$--$8$ & $10$--$21$ & $S_3$--$S_4$ & Near-complete coverage; institutional lock-in\\
    \textbf{Universal} & $\geq 9$ & $\geq 28$ & $S_4$--$S_6$ & Asymptotic saturation; $\CECstar$ negligible\\
    \bottomrule
  \end{tabular}
\end{table}

\subsection{The Law of Topological Acceleration}

The relationship derived in Proposition~\ref{prop:cec_reduction} implies what we term the \textbf{Law of Topological Acceleration}:

\begin{quote}
  \emph{Each additional expert reduces the critical CEC threshold $\CECstar$ by a factor that scales as the incremental Betti number $\Delta\beta_1(M) = M-1$. The adoption convergence time $\tau_{\text{adopt}}$, being proportional to $\CECstar$, inherits the same scaling:}
  \begin{equation}\label{eq:tau_adopt}
    \tau_{\text{adopt}}(M) \propto \frac{1}{1 + \alpha \cdot \beta_1(M)}.
  \end{equation}
\end{quote}

This law has immediate practical implications: to accelerate adoption convergence, increase the expert pool size. The marginal benefit of adding the $(M+1)$-th expert is $\Delta\tau \propto \alpha(M-1)$, growing linearly with $M$. Large expert pools converge to lock-in exponentially faster than small ones.

% ══════════════════════════════════════════════════════════════════════
% 6. IMPLICATIONS FOR PROTOCOL DESIGN
% ══════════════════════════════════════════════════════════════════════
\section{Implications for Protocol Design}

\subsection{Minimum Expert Count for Robust Adoption}

The analysis yields a clear lower bound:
\begin{quote}
  \textbf{Design Principle 1 (Topological Minimum).} \emph{Any multi-expert audit protocol intended to achieve non-proliferation equilibrium ($S_3$) must operate with at least $M \geq 4$ independent experts, corresponding to $\beta_1 \geq 3$.}
\end{quote}

$M=3$ ($\beta_1=1$) provides the \emph{capability} of detecting irreducible inconsistency, but with only one harmonic channel, the statistical power is insufficient to drive $\CECstar$ low enough for practical lock-in time scales---unless the topological coupling $\alpha$ is unexpectedly large. $M=4$ adds two additional channels, tripling the effective detection capacity and providing the triangulation necessary for reliable consensus.

\subsection{Optimal Auditor Pool Size}

The optimal pool size involves a trade-off:
\begin{itemize}
  \item \textbf{Larger $M$:} Reduces $\CECstar$ and accelerates adoption convergence (topological benefit).
  \item \textbf{Larger $M$:} Increases coordination cost, communication overhead, and the risk of expert collusion (economic cost).
\end{itemize}

Let $C(M) = c_0 + c_1 M$ be the per-audit operating cost for $M$ experts, and let $V(\CECstar(M))$ be the value of achieving lock-in at CEC size $\CECstar(M)$. The net social value as a function of $M$ is
\begin{equation}
  W(M) = V(\CECstar(M)) - C(M) \cdot \tau_{\text{adopt}}(M).
\end{equation}

Using $\CECstar(M) \propto (1 + \alpha\beta_1)^{-1}$ and $\tau_{\text{adopt}} \propto \CECstar$, we obtain
\begin{equation}
  W(M) \approx V_0 \cdot \frac{\alpha\beta_1(M)}{1+\alpha\beta_1(M)} - \frac{C(M)}{1+\alpha\beta_1(M)}.
\end{equation}

With $\beta_1(M) \sim M^2/2$, the first term saturates as $V_0 \cdot \alpha M^2/(2 + \alpha M^2)$, while the cost term decays as $2C(M)/(2+\alpha M^2)$. The optimum occurs at intermediate $M$ where the marginal topological benefit equals the marginal cost.

\begin{proposition}[Optimal Expert Count]\label{prop:optimal_M}
  For linear cost $C(M) = c_0 + c_1 M$, the optimal expert count $M^*$ satisfies
  \begin{equation}
    M^* \approx \sqrt{\frac{2V_0}{c_1}} \cdot \alpha^{-1/2},
  \end{equation}
  in the regime where $\alpha M^2 \gg 2$ and $c_0$ is negligible.
\end{proposition}

For typical parameter values ($V_0/c_1 \sim 10^2$, $\alpha \sim 0.1$), this yields $M^* \approx 14$, consistent with the intuition that audit pools of 10--15 experts provide near-optimal topological benefits before coordination costs dominate.

\subsection{Topological Redundancy and Fault Tolerance}

A further implication concerns robustness to expert failure or defection. If one of $M$ experts becomes compromised or unavailable, the effective Betti number drops not to $\beta_1(M-1)$ but to a value dependent on \emph{which} expert is lost---because the harmonic modes are global structures, not local to any single vertex. The worst-case reduction is
\begin{equation}
  \beta_1^{\text{(after loss)}} \geq \beta_1(M) - (M-2) = \frac{(M-1)(M-2)}{2} - (M-2) = \frac{(M-2)(M-3)}{2}.
\end{equation}

This implies that the system retains $\beta_1(M-1) = (M-2)(M-3)/2$ independent modes after the loss of \emph{any} single expert---a form of topological fault tolerance. The loss of one expert among $M$ reduces the inconsistency detection space by exactly $M-2$ dimensions, preserving the remaining structure.

\begin{corollary}[Topological Fault Tolerance]\label{cor:fault_tolerance}
  The system with $M$ experts retains $\beta_1(M-1)$ independent harmonic modes after the loss of any single expert. The fractional loss is
  \begin{equation}
    \frac{\Delta\beta_1}{\beta_1} = \frac{M-2}{\beta_1(M)} = \frac{2}{M-1},
  \end{equation}
  decaying as $2/M$ for large $M$.
\end{corollary}

For $M=4$, losing one expert reduces $\beta_1$ from 3 to 0 (a 100\% loss), confirming that $M=4$ is the \emph{absolute minimum} for topological fault tolerance. For $M=10$, losing one expert reduces $\beta_1$ from 36 to 28 (a 22\% loss), a manageable degradation.

% ══════════════════════════════════════════════════════════════════════
% 7. CONJECTURE
% ══════════════════════════════════════════════════════════════════════
\section{Main Conjecture}

The analysis has established a scaling relationship between $\beta_1$ and $\CECstar$ under the specific functional form~\eqref{eq:theta_enhanced}. But the precise functional dependence of accuracy on $\beta_1$ is not uniquely determined by the theory as it currently stands. We therefore formulate:

\begin{conjecture}[Topological Adoption Conjecture]\label{conj:main}
  \textbf{Universal scaling.} For any multi-expert audit system satisfying the \SCX\ assumptions (independent training data, state-homogeneous expert quality, bounded label noise), the critical CEC threshold scales as
  \begin{equation}\label{eq:conjecture}
    \CECstar(M) \propto \frac{1}{\beta_1(M)} \quad \text{as } M \to \infty,
  \end{equation}
  independent of the precise functional form of the accuracy curve $\theta(|\CEC|)$ and the topological coupling $\alpha$. The scaling exponent is $-1$ in the Betti number.
\end{conjecture}

\textbf{Motivation.} The $\beta_1$ harmonic modes are mutually orthogonal in the Hodge inner product. In the large-$M$ limit, the expert graph approaches a complete graph $K_M$, whose harmonic space is the full space of 1-cycles modulo boundaries. Each harmonic mode contributes an independent degree of freedom to the inconsistency detection, and by the central limit theorem, the combined detection power scales linearly with the number of independent modes. Since $\CECstar$ is inversely proportional to the per-audit detection efficiency (by the exponential accuracy model), and detection efficiency scales with the effective number of independent tests, we expect $\CECstar \propto 1/\beta_1$ asymptotically.

\textbf{The $\alpha \to 1$ limit.} The conjecture is most naturally interpreted in the limit where the topological coupling $\alpha \to 1$, i.e., where each harmonic mode contributes its maximum possible detection efficiency. In this limit, $\CECstar(M) = \CECstar_0/(1 + \beta_1) \sim \CECstar_0/\beta_1$ for large $M$, recovering the conjectured scaling.

\textbf{Testable prediction.} The conjecture predicts that systems with $M=10$ experts ($\beta_1=36$) achieve lock-in approximately $36\times$ faster (in CEC accumulation) than the $M=2$ baseline ($\beta_1=0$, but effectively $\beta_1 = 1$ as the normalization baseline). This is a quantitative prediction amenable to simulation or empirical testing once multi-expert audit systems of sufficient scale are deployed.

% ══════════════════════════════════════════════════════════════════════
% 8. CONCLUSION
% ══════════════════════════════════════════════════════════════════════
\section{Conclusion and Open Problems}

\subsection{Summary of Contributions}

We have established the following chain of reasoning:

\begin{enumerate}[nosep]
  \item The \SCX\ expert graph's first Betti number $\beta_1 = (M-1)(M-2)/2$ is a topological invariant counting independent harmonic modes---incommensurable flat gauge connections that represent irreducible inconsistencies among experts.
  \item $\beta_1$ exhibits discontinuous jumps at each increment of $M$: $0 \to 1 \to 3 \to 6 \to 10 \to \cdots$
  \item These jumps enhance multi-expert detection efficiency by providing independent inconsistency channels, multiplicatively reducing the critical CEC threshold $\CECstar$ required for adoption lock-in.
  \item The relationship $\CECstar(M) = \CECstar_0/(1 + \alpha\beta_1(M))$ yields a topological acceleration law: $\tau_{\text{adopt}} \propto 1/\beta_1$ for large $M$.
  \item The $\beta_1$ jumps map structurally onto the adoption phase transitions $S_1 \to S_2 \to \cdots \to S_6$, with $M=4$ ($\beta_1=3$) identified as the minimal topologically robust expert count.
\end{enumerate}

\subsection{Open Problems}

\begin{problem}[Empirical Determination of $\alpha$]
  The topological coupling constant $\alpha$---the per-mode detection efficiency gain---is not determined from first principles. It must be measured empirically or bounded through simulation. What is the range of $\alpha$ for realistic expert ensembles?
\end{problem}

\begin{problem}[Topological Phase Transition Rigor]
  The mapping from $\beta_1$ jumps to adoption transitions is proposed on structural grounds. Can it be proved as a theorem, i.e., that $\CECstar(M)$ being below a certain threshold is a \emph{necessary and sufficient} condition for the $S_i \to S_{i+1}$ transition?
\end{problem}

\begin{problem}[Non-Abelian Generalization]
  The current analysis uses the Abelian (translation) gauge group $\R^d$. For $O(d)$ gauge (rotation), the harmonic moduli space dimension is $\beta_1 \cdot d(d-1)/2$. Does the rotation gauge introduce additional detection channels that further accelerate adoption? If so, the effective $\beta_1$ for $O(d)$ may be larger than the combinatorial $\beta_1$.
\end{problem}

\begin{problem}[Dynamic $\beta_1$ in Spring Self-Evolution]
  The \Spring\ self-evolution loop may effectively increase $\beta_1$ over time as the cumulative evidentiary corpus grows, by revealing latent harmonic modes that were not initially detectable. Does $\beta_1$ itself evolve dynamically, and if so, does this create a positive feedback loop with adoption dynamics?
\end{problem}

\begin{problem}[Verification of the Conjecture]
  The Main Conjecture ($\CECstar \propto 1/\beta_1$) predicts a specific asymptotic scaling. Can this be verified through large-scale simulation of multi-expert audit systems with variable $M$?
\end{problem}

\subsection{Final Remark}

The unification of gauge-theoretic topology and game-theoretic adoption dynamics presented here is, at present, a \emph{prototype}---a structural proposal supported by mathematical analogy and scaling arguments rather than rigorous theorem-proving. We regard this as appropriate for a first exploration of the connection. The topology of the \SCX\ graph is real, mathematically precise, and carries genuine information about the system's capacity to represent and detect inconsistency. The adoption dynamics are real, game-theoretically formalized, and exhibit phase transitions that quantum field theorists would recognize as analogs of symmetry breaking. The structural correspondence between them is too systematic to be coincidental; whether it rises to the level of a theorem or remains a fruitful heuristic is the question this prototype paper is designed to provoke.

% ══════════════════════════════════════════════════════════════════════
% ACKNOWLEDGMENTS
% ══════════════════════════════════════════════════════════════════════
\section*{Acknowledgments}
The authors acknowledge the foundational work of the \SCX\ gauge theory team on the discrete Hodge decomposition and Dijkgraaf-Witten TQFT framework, and the \Yajie\ Protocol team on the game-theoretic adoption model and lock-in trajectory analysis. This paper is a synthesis across these two lines of development and would not have been possible without both.

% ══════════════════════════════════════════════════════════════════════
% BIBLIOGRAPHY
% ══════════════════════════════════════════════════════════════════════
\begin{thebibliography}{99}

\bibitem{scx_gauge_formalized}
\SCX\ Research Group.
\emph{SCX Gauge Theory Formalized: $O(d)$ Lattice Gauge, Dijkgraaf-Witten TQFT, and Information-Geometric Bulk-Boundary Correspondence.}
SCX Papers, v2.0, 2026.

\bibitem{scx_fiber_bundle}
\SCX\ Research Group.
\emph{Discrete Geometry of PES Misalignment: A Graph Hodge Formalization of SCX Gauge Theory.}
SCX Papers, 2026.

\bibitem{scx_theory}
\SCX\ Research Group.
\emph{The SCX Theory: Theorems, Concepts, and Assumptions.}
SCX Papers, 2026.

\bibitem{yajie_protocol}
Yajie Research Group.
\emph{The Yajie Protocol: Technology Lock-in, Audit Sovereignty, and the Non-Proliferation Logic of Data Quality Assessment.}
Journal of Data Governance and Information Systems (Submitted), June 2026.

\bibitem{protocol_governance}
\SCX\ Research Group.
\emph{Protocol Governance: Game-Theoretic Foundations of the SCX Maintainer Rotation Mechanism.}
SCX Papers, 2026.

\bibitem{gauge_domain_formalization}
\SCX\ Research Group.
\emph{Gauge Domain Formalization: Rigorous Theorem Statements for Revived Domains.}
SCX Supplementary Analysis, 2026.

\bibitem{gauge_inventory}
\SCX\ Research Group.
\emph{SCX Gauge Theory Complete Inventory.}
SCX Internal Documentation, July 2026.

\bibitem{gametheory_inventory}
\SCX\ Research Group.
\emph{SCX Game Theory \& Economics Complete Inventory.}
SCX Internal Documentation, July 2026.

\end{thebibliography}

\end{document}
